\section{Related Work}
\label{sec:related}

\parabf{Inference serving.} There has been plenty of work on inference serving recently. They range from general-purpose production-grade systems like TorchServe~\cite{torchserve} and NVIDIA Triton~\cite{nvidiatriton} to systems optimized specifically for Transformer-based LLMs~\cite{li2023alpaserve, yu2022orca, agrawal2023sarathi, wu2023fast, fang2021turbotransformers,fastertransformer,zhou2022pets,su2023hotgpt}. 
Among them, Orca~\cite{yu2022orca} introduces continuous batching to increase throughput. vLLM~\cite{vllm} proposes paged-attention for fine-grained KV cache management. SARATHI~\cite{agrawal2023sarathi} suggests a chunked-prefill approach, splitting a prefill request into chunks and piggybacking decoding requests to improve hardware utilization. FastServe~\cite{wu2023fast} implements iteration-level preemptive scheduling to mitigate the queuing delay caused by long jobs. However, they all employ a colocation approach for prefill and decoding processing, thus leading to severe interference. There are also concurrent works such as Splitwise~\cite{patel2023splitwise}, TetriInfer~\cite{hu2024inference} and DéjàVu~\cite{strati2024dejavu} which adopt similar disaggregation idea to optimize LLM inference, further confirming the effectiveness of this method. Differently, \sys emphasizes the goodput optimization scenario more and takes a closer look at the aspect of network bandwidth.


\parabf{Goodput-optimized systems.} Optimizing goodput is a hot topic in DL applications. Pollux~\cite{pollux} improves scheduling performance in DL clusters by dynamically adjusting resources for jobs to increase cluster-wide goodput. Sia~\cite{jayaram2023sia} introduces a heterogeneous-aware scheduling approach that can efficiently match cluster resources to elastic resource-adaptive jobs. Clockwork~\cite{clockwork} and Shepherd~\cite{zhangshepherd} provide latency-aware scheduling and preemption to improve the serving goodput, but they only target traditional small models. AlpaServe~\cite{li2023alpaserve} focuses on LLMs, employing model parallelism to statistically multiplex the GPU execution thus improving the resource utilization. However, it only targets the non-autoregressive generation. \sys is the first work to optimize the goodput for autoregressive LLM inference.


\parabf{Resource disaggregation.} Resource disaggregated systems~\cite{shan2018legoos,mira,CXL} decouple the hardware resources from the traditional monolithic server infrastructure and separate them into resource pools to manage independently. It allows for more flexible, efficient, and scalable deployment and increases resource utilization. Many applications benefit from a truly disaggregated data center with high-speed network bandwidth and heterogenous hardware support ~\cite{makeitreal,audibert2022case, distmind}. \sys shares the concept by disaggregating its system components, allowing for independent resource scaling and management.


\parabf{Model parallelism for training.} \sys is orthogonal to the large body of work on model parallelism in training~\cite{alpa,shoeybi2020megatronlm,huang2019gpipe,rajbhandari2020zero,pipedream}. As described in \S\ref{subsec:practical_problems}, inference-serving workloads have unique characteristics not found in training settings. Where these systems do intersect with \sys, is in their methods for implementing model parallelism along various dimensions. \sys can integrate new parallelism optimizations into its placement searching algorithm.
